\begin{claim}{D.60}{Proved}{The exact torsion quotient of a negative torus extension}
Use D.59's notation
\[
 E=F\rtimes\langle v\rangle,\quad F=\widehat F(x,y),
 \quad c=[x,y],\quad v^2=1,\quad
 \alpha=\operatorname{Inn}(v)|_F,\quad M=F^{\rm ab},
\]
and choose an actual transporter $g$ with
$gvg^{-1}=cv$.
For a subset of a profinite group, double angle
brackets below mean its closed normal closure in
the specified ambient group. There are exact
continuous quotient identifications
\[
 \begin{split}
 E/\langle\!\langle v\rangle\!\rangle_E
 &\ \xrightarrow{\ \cong\ }\ M/(\alpha_*-I)M,\\
 E/\langle\!\langle\operatorname{Tor}(E)\rangle\!\rangle_E
 &\ \xrightarrow{\ \cong\ }\ \widehat\Z ,
 \end{split}                                             \tag{D.60a}
\]
where $\operatorname{Tor}(E)$ is the set of
finite-order elements. More explicitly,
\[
\begin{array}{c|c|c}
 A_2\bmod2&
 E/\langle\!\langle v\rangle\!\rangle_E&
 \chi(g)\\ \hline
 I&\widehat\Z\times C_2&1\\
 \ne I&\widehat\Z&2 .
\end{array}                                              \tag{D.60b}
\]
Here $\chi:E\to\widehat\Z$ is the second quotient
in (D.60a), with the normalization specified in
the proof. Its restriction to $F$ is onto.
It has a continuous section whose image lies
in $F$, so
\[
 E=\langle\!\langle\operatorname{Tor}(E)\rangle\!\rangle_E
                         \rtimes s(\widehat\Z).          \tag{D.60c}
\]
In the first row, let $w$ represent the other
full section class. One can choose that section
$s$ inside $C_F(w)$, and the resulting map
\[
 \widehat\Z\times C_2\longrightarrow E,\qquad
            (t,\epsilon)\longmapsto s(t)w^\epsilon       \tag{D.60d}
\]
is a continuous embedding. These statements specify
the quotient, splitting, and embedded subgroup;
they give no further free-product presentation
of $E$.
\end{claim}

\begin{proof}
\emph{The two quotient maps before any module decomposition.}
Write $R=\langle\!\langle v\rangle\!\rangle_E$ and
$Q=E/R$. The map $F\to Q$ is onto, since $E$
is generated by $F$ and $v$.
D.58's actual transporter gives
\[
                  c=[g,v]=g\alpha(g)^{-1}.               \tag{D.60e}
\]
Consequently $c$ dies in $Q$.
The profinite quotient
$F/\langle\!\langle[x,y]\rangle\!\rangle_F$
is exactly $\widehat\Z^2=M$: its two generating
images commute, while the abelianization map
to $\widehat\Z^2$ gives the inverse universal map.
Thus the specified map $F\to Q$ factors through
$M$.
Since $v$ dies, $\alpha(f)$ and $f$ have the same
image for every $f\in F$. This supplies a
continuous surjection
\[
                 M/(\alpha_*-I)M\longrightarrow Q.       \tag{D.60f}
\]
The image $(\alpha_*-I)M$ is compact and hence
closed.

In the other direction, the formula
\[
 E\longrightarrow M/(\alpha_*-I)M,\qquad
           f v^\epsilon\longmapsto [f]\bmod(\alpha_*-I)M
                                                               \tag{D.60g}
\]
is a continuous homomorphism. Indeed the product
in the specified semidirect extension is
$f\alpha^\epsilon(h)v^{\epsilon+\eta}$, and
$\alpha_*$ acts trivially on this quotient.
Formula (D.60g) kills $v$, so it factors through
$Q$. Its composite with (D.60f) is the identity
on the images of $F$; the reverse composite has
the same property. Those images generate both
quotients, proving the first isomorphism in
(D.60a).
Equivalently, all of its full-group relations
can be retained as
\[
 Q=F/\langle\!\langle
       c,\ \alpha(f)f^{-1}\ (f\in F)
                         \rangle\!\rangle_F .
                                                               \tag{D.60h}
\]
The relation $c$ is already implied by the
displayed relation with $f=g$, through (D.60e).

\emph{Explicit coordinates for both integral module cases.}
Put $R_0=\widehat\Z$.
D.59 proves that $e_+=[g]$ generates the entire
fixed line in $M$ and that this line is a direct
summand. Choose a complementary basis vector
$e$ in the \emph{homology module}, so
$M=R_0 e_+\oplus R_0 e$.
No free generating pair of $F$ is changed.
Since the global determinant of $\alpha_*$ is
$-1$ by D.43, its matrix in these coordinates is
\[
 A=\begin{pmatrix}1&a\\0&-1\end{pmatrix},\qquad
  (A-I)M=R_0(a,-2)                                      \tag{D.60i}
\]
for a specified $a\in R_0$.
The two cases of (D.60b) are exactly whether
$a\in2R_0$ or $a_2\in\Z_2^\times$.

In the first case write $a=2b$.
The following map is onto:
\[
 R_0^2\longrightarrow R_0\times C_2,\qquad
                   (r,t)\longmapsto(r+bt,\ t\bmod2).
                                                               \tag{D.60j}
\]
Its kernel is exactly $R_0(a,-2)$.
Indeed $t=2z$ and $r+bt=0$ are equivalent to
$(r,t)=(-az,2z)$.
Conversely each multiple of $(a,-2)$ has zero
image. The map is onto by choosing $t=0$ or $1$
for the prescribed second coordinate and then
choosing $r$ for the first.
It therefore gives the first-row quotient
homeomorphism, since a continuous bijection
between compact Hausdorff spaces is a
homeomorphism.

In the second case the map
\[
 R_0^2\longrightarrow R_0,\qquad
                         (r,t)\longmapsto2r+at          \tag{D.60k}
\]
is onto.
At every odd prime its coefficient $2$ is a
unit; at the prime two its coefficient $a_2$
is a unit. These coordinatewise solutions
give an element of $R_0^2$ for every prescribed
value. Its kernel is again exactly
$R_0(a,-2)$:
if $2r+at=0$, the dyadic unit $a_2$ forces
$t_2\in2\Z_2$, hence $t\in2R_0$.
Write $t=-2z$. The equation becomes
$2(r-az)=0$, and $R_0$ has no additive
two-torsion, so $r=az$.
The converse follows by substitution.
This gives the second-row quotient
homeomorphism.

For $[f]=r e_++t e$, define $\chi(fv^\epsilon)$
to be $r+bt$ in the first case and $2r+at$
in the second. Formulas (D.60g), (D.60j), and
(D.60k) prove that these are continuous
homomorphisms. They give exactly the values
$\chi(g)=1$ and $2$, respectively, and
$\chi|_F$ is onto in both cases.

\emph{The second section and the precise finite factor.}
The free profinite group $F$ is torsion-free.
Every nontrivial finite-order element of $E$
therefore has order two and generates a section:
its finite cyclic subgroup injects into
$E/F=C_2$.
D.58 gives one full $F$-conjugacy class in the
second row and exactly two in the first.

In the first row, the dyadic coordinate
character $[f]=r e_++t e\mapsto t_2$ is a
primitive anti-invariant character
$\varphi_-:F\to\Z_2$, since (D.60i) sends
$t$ to $-t$.
Its mod-two extension
\[
 \beta:E\longrightarrow\F_2,\qquad
                    \beta(fv^\epsilon)=t_2\bmod2        \tag{D.60l}
\]
is a homomorphism and vanishes on $v$.
It is exactly the D.58 section-distinguishing
character normalized to vanish on this cusp
reflection.
To make that normalization explicit, if D.58
was written using another initial section $s$,
its character $b$ has the same restriction
to $F$ as (D.60l), and
\[
                         \beta=b+b(v)u,                 \tag{D.60m}
\]
where $u:E\to C_2$ is the original quotient
character. Primitive choices of the
anti-invariant integral character differ by
a dyadic unit and have the same mod-two
restriction. D.58 proves that the two section
classes have opposite $b$-values.
Equation (D.60m) thus gives
$\beta(v)=0$ and $\beta(w)=1$ for every
representative $w$ of the other class.

The image of $w$ in the quotient (D.60j)
has order dividing two.
The additive group $R_0$ is torsion-free,
so this image has first coordinate zero.
Equation (D.60l) is its second coordinate
and has value one. Thus that image is
exactly $(0,1)$.
The closed normal subgroup generated by
all torsion in the first row is consequently
the inverse image of $\{0\}\times C_2$ under
$E\to Q$.
In the second row all torsion is conjugate
to $v$, so its closed normal closure is
already $R$.
This proves in both cases
\[
 \ker\chi=
       \langle\!\langle\operatorname{Tor}(E)\rangle\!\rangle_E
                                                               \tag{D.60n}
\]
and proves the second isomorphism in (D.60a).
It also identifies the exact relation between
the two quotient maps, including the entire
finite factor in the first row.

\emph{Actual splittings and the embedded commuting factor.}
Choose $h\in F$ with $\chi(h)=1$, possible
because $\chi|_F$ is onto.
The continuous profinite power map
\[
                 s:R_0\longrightarrow E,\qquad t\longmapsto h^t
                                                               \tag{D.60o}
\]
satisfies $\chi(s(t))=t$.
It is therefore injective, and compactness
identifies its image with $R_0$ topologically.
Every $z\in E$ has the unique factorization
\[
                  z=\bigl(zs(\chi(z))^{-1}\bigr)s(\chi(z)),
 \qquad zs(\chi(z))^{-1}\in\ker\chi .
\]
This proves (D.60c), with the given maps.

In the first row, D.59 proves that
$D=C_F(w)\to M^+$ is onto.
Choose $h\in D$ with $[h]=[g]$.
Then $\chi(h)=1$, so the same construction
gives a section $s$ inside $D$.
The formula (D.60d) is a homomorphism
because $h$ commutes with $w$.
If its value is the identity, application
of $\chi$ gives $t=0$, and then
$w^\epsilon=1$ gives $\epsilon=0$.
It is a continuous embedding by compactness.
This establishes every stated splitting and
subgroup without identifying their generated
subgroup with the whole extension.

The executable
\texttt{code/part\_c/torus\_torsion\_quotient\_d60.py}
checks the displayed quotient maps and their
exact kernels in finite coefficient rings.
The full quotient identifications follow from
the two universal maps and the actual
transporter relation above.
\end{proof}
